% --- Archon physics formalization source begin ---
% archon:physics
% archon:covers PhyXMiniProblems/problem_phyx_mini_0735.lean
% archon:source-report reports/phyx_mini/problem_phyx_mini_0735.source.json

\chapter{Physics problem phyx\_mini\_0735}
\label{ch:PhyXMiniProblems_problem_phyx_mini_0735}

\paragraph{Problem source.}
\#\# Physical scenario

Rock faults are ruptures along which opposite faces of rock have slid past each other. In figure, points \$A\$ and \$B\$ coincided before the rock in the foreground slid down to the right. The net displacement \$\textbackslash{}vec\{AB\}\$ is along the plane of the fault. The horizontal component of \$\textbackslash{}vec\{AB\}\$ is the strike-slip \$AC\$. The component of \$\textbackslash{}vec\{AB\}\$ that is directed down the plane of the fault is the dip-slip \$AD\$.

\#\# Question

What is the magnitude of the net displacement \$\textbackslash{}vec\{AB\}\$ if the strike-slip is \$22.0\textbackslash{} m\$ and the dip-slip is \$17.0\textbackslash{} m\$?

\#\# Answer choices

- A: A: 23.7 m

- B: B: 25.4 m

- C: C: 27.8 m

- D: D: 29.6 m

\#\# Figure caption (auxiliary; use the image as primary evidence)

The image depicts a diagram of a fault plane, commonly used in geology to illustrate fault movements. It shows two blocks with a fault plane between them. Here's a detailed description:

1. **Blocks**: 
   - The diagram features two blocks, each representing sections of the Earth's crust. They are textured to resemble rock surfaces.
   
2. **Fault Plane**: 
   - A central diagonal surface between the two blocks represents the fault plane. It is labeled accordingly.

3. **Arrows and Lines**: 
   - An orange arrow points from point A to point B, indicating movement along the fault plane.
   - A dashed line between points C and D represents the horizontal line on the fault plane.

4. **Labels**: 
   - "Strike-slip" is labeled along the horizontal dimension, indicating lateral movement along the plane.
   - "Dip-slip" is labeled along the vertical dimension, indicating vertical movement.
   
5. **Angles and Bearings**: 
   - There is a red angle marked at point A, where the meeting lines of strike-slip and dip-slip are highlighted. 
   - The Greek letter φ (phi) denotes the angle between the dashed line and the vertical, representing the dip angle of the fault plane.

6. **Points**: 
   - Points A, B, C, and D mark key intersections and angles within the diagram to illustrate geometric relationships.

The diagram effectively shows the components and mechanics involved in fault plane movements, useful for understanding geological shifts.

\#\# Dataset metadata

Kinematics; Spatial Relation Reasoning, Multi-Formula Reasoning

\paragraph{Recorded answer/context.}
C: C: 27.8 m

\paragraph{Figure/image path.}
/root/proposal\_for\_physic/science-mango/phyx\_mini\_run/phyx\_data/test\_image/735.png

\paragraph{Formalization target.}
create a compiling Lean file with sorry bodies at `PhyXMiniProblems/problem\_phyx\_mini\_0735.lean`. The Lean declarations must preserve the physical quantities, dimensions or dimensional roles, figure labels, governing-law hypotheses, and final relation expressed by this problem.
Use Mathlib/Physlib names found through LeanExplore where available. If a physics API is missing, introduce faithful local abstractions rather than scalar placeholder aliases.

\begin{theorem}[Physics formalization target]
\label{thm:physics:phyx_mini_0735:target}
\lean{PhyXMiniProblems.ProblemPhyXMini0735.netDisplacementMagnitude}
\uses{def:physics:phyx-mini-0735:phyxminiproblems-problemphyxmini0735-displacementmagnitudeinmeters, def:physics:phyx-mini-0735:phyxminiproblems-problemphyxmini0735-faultplanefigure, def:physics:phyx-mini-0735:phyxminiproblems-problemphyxmini0735-matchesprimaryfigure, def:physics:phyx-mini-0735:phyxminiproblems-problemphyxmini0735-faultcomponentlaws, def:physics:phyx-mini-0735:phyxminiproblems-problemphyxmini0735-matchesproblemdata}
The assigned autoformalize agent should translate this physics problem into Lean declarations in the covered file, with theorem and lemma proof bodies written as `by sorry`.
\end{theorem}
\begin{proof}
This is an autoformalization task, not a proof task. Produce statements that can later be proved by the physics prover without weakening the physical model.
\end{proof}

% --- Archon declaration topology begin ---
\section{Declaration topology}
\begin{definition}[Fault Plane Vector]
  \label{def:physics:phyx-mini-0735:phyxminiproblems-problemphyxmini0735-faultplanevector}
  \lean{PhyXMiniProblems.ProblemPhyXMini0735.FaultPlaneVector}
  A real two-dimensional coordinate vector intrinsic to the fault plane
\end{definition}
\begin{proof}
  This is the declared model definition of Fault Plane Vector.
\end{proof}
\begin{definition}[Displacement Vector]
  \label{def:physics:phyx-mini-0735:phyxminiproblems-problemphyxmini0735-displacementvector}
  \lean{PhyXMiniProblems.ProblemPhyXMini0735.DisplacementVector}
  \uses{def:physics:phyx-mini-0735:phyxminiproblems-problemphyxmini0735-faultplanevector}
  A unit-independent physical displacement vector of length dimension
\end{definition}
\begin{proof}
  This is the declared model definition of Displacement Vector.
\end{proof}
\begin{definition}[displacement Readout]
  \label{def:physics:phyx-mini-0735:phyxminiproblems-problemphyxmini0735-displacementreadout}
  \lean{PhyXMiniProblems.ProblemPhyXMini0735.displacementReadout}
  \uses{def:physics:phyx-mini-0735:phyxminiproblems-problemphyxmini0735-faultplanevector, def:physics:phyx-mini-0735:phyxminiproblems-problemphyxmini0735-displacementvector}
  Read a physical displacement vector in a selected length unit
\end{definition}
\begin{proof}
  This is the declared model definition of displacement Readout.
\end{proof}
\begin{definition}[displacement In Meters]
  \label{def:physics:phyx-mini-0735:phyxminiproblems-problemphyxmini0735-displacementinmeters}
  \lean{PhyXMiniProblems.ProblemPhyXMini0735.displacementInMeters}
  \uses{def:physics:phyx-mini-0735:phyxminiproblems-problemphyxmini0735-faultplanevector, def:physics:phyx-mini-0735:phyxminiproblems-problemphyxmini0735-displacementvector, def:physics:phyx-mini-0735:phyxminiproblems-problemphyxmini0735-displacementreadout}
  Read a physical displacement vector in metres
\end{definition}
\begin{proof}
  This is the declared model definition of displacement In Meters.
\end{proof}
\begin{definition}[displacement Magnitude Readout]
  \label{def:physics:phyx-mini-0735:phyxminiproblems-problemphyxmini0735-displacementmagnitudereadout}
  \lean{PhyXMiniProblems.ProblemPhyXMini0735.displacementMagnitudeReadout}
  \uses{def:physics:phyx-mini-0735:phyxminiproblems-problemphyxmini0735-displacementvector, def:physics:phyx-mini-0735:phyxminiproblems-problemphyxmini0735-displacementreadout}
  Read the magnitude of a physical displacement in a selected length unit
\end{definition}
\begin{proof}
  This is the declared model definition of displacement Magnitude Readout.
\end{proof}
\begin{definition}[displacement Magnitude In Meters]
  \label{def:physics:phyx-mini-0735:phyxminiproblems-problemphyxmini0735-displacementmagnitudeinmeters}
  \lean{PhyXMiniProblems.ProblemPhyXMini0735.displacementMagnitudeInMeters}
  \uses{def:physics:phyx-mini-0735:phyxminiproblems-problemphyxmini0735-displacementvector, def:physics:phyx-mini-0735:phyxminiproblems-problemphyxmini0735-displacementmagnitudereadout}
  Read the magnitude of a physical displacement in metres
\end{definition}
\begin{proof}
  This is the declared model definition of displacement Magnitude In Meters.
\end{proof}
\begin{definition}[Fault Point]
  \label{def:physics:phyx-mini-0735:phyxminiproblems-problemphyxmini0735-faultpoint}
  \lean{PhyXMiniProblems.ProblemPhyXMini0735.FaultPoint}
  The four point labels printed on the supplied fault-plane diagram
\end{definition}
\begin{proof}
  This is the declared model definition of Fault Point.
\end{proof}
\begin{definition}[Fault Segment]
  \label{def:physics:phyx-mini-0735:phyxminiproblems-problemphyxmini0735-faultsegment}
  \lean{PhyXMiniProblems.ProblemPhyXMini0735.FaultSegment}
  The labeled segments used in the parallelogram decomposition
\end{definition}
\begin{proof}
  This is the declared model definition of Fault Segment.
\end{proof}
\begin{definition}[segment Start]
  \label{def:physics:phyx-mini-0735:phyxminiproblems-problemphyxmini0735-segmentstart}
  \lean{PhyXMiniProblems.ProblemPhyXMini0735.segmentStart}
  \uses{def:physics:phyx-mini-0735:phyxminiproblems-problemphyxmini0735-faultpoint, def:physics:phyx-mini-0735:phyxminiproblems-problemphyxmini0735-faultsegment}
  The initial endpoint associated with each labeled segment
\end{definition}
\begin{proof}
  This is the declared model definition of segment Start.
\end{proof}
\begin{definition}[segment End]
  \label{def:physics:phyx-mini-0735:phyxminiproblems-problemphyxmini0735-segmentend}
  \lean{PhyXMiniProblems.ProblemPhyXMini0735.segmentEnd}
  \uses{def:physics:phyx-mini-0735:phyxminiproblems-problemphyxmini0735-faultpoint, def:physics:phyx-mini-0735:phyxminiproblems-problemphyxmini0735-faultsegment}
  The terminal endpoint associated with each labeled segment
\end{definition}
\begin{proof}
  This is the declared model definition of segment End.
\end{proof}
\begin{definition}[Segment Role]
  \label{def:physics:phyx-mini-0735:phyxminiproblems-problemphyxmini0735-segmentrole}
  \lean{PhyXMiniProblems.ProblemPhyXMini0735.SegmentRole}
  Kinematic roles assigned to lines in the supplied figure
\end{definition}
\begin{proof}
  This is the declared model definition of Segment Role.
\end{proof}
\begin{definition}[Segment Style]
  \label{def:physics:phyx-mini-0735:phyxminiproblems-problemphyxmini0735-segmentstyle}
  \lean{PhyXMiniProblems.ProblemPhyXMini0735.SegmentStyle}
  Drawing styles that distinguish the figure's arrow and construction lines
\end{definition}
\begin{proof}
  This is the declared model definition of Segment Style.
\end{proof}
\begin{definition}[Fault Plane Figure]
  \label{def:physics:phyx-mini-0735:phyxminiproblems-problemphyxmini0735-faultplanefigure}
  \lean{PhyXMiniProblems.ProblemPhyXMini0735.FaultPlaneFigure}
  \uses{def:physics:phyx-mini-0735:phyxminiproblems-problemphyxmini0735-displacementvector, def:physics:phyx-mini-0735:phyxminiproblems-problemphyxmini0735-faultpoint, def:physics:phyx-mini-0735:phyxminiproblems-problemphyxmini0735-faultsegment, def:physics:phyx-mini-0735:phyxminiproblems-problemphyxmini0735-segmentrole, def:physics:phyx-mini-0735:phyxminiproblems-problemphyxmini0735-segmentstyle}
  The physical vectors and auxiliary annotations displayed in image 735. The dip angle phi is dimensionless and is represented by its radian readout
\end{definition}
\begin{proof}
  This is the declared model definition of Fault Plane Figure.
\end{proof}
\begin{definition}[Matches Primary Figure]
  \label{def:physics:phyx-mini-0735:phyxminiproblems-problemphyxmini0735-matchesprimaryfigure}
  \lean{PhyXMiniProblems.ProblemPhyXMini0735.MatchesPrimaryFigure}
  \uses{def:physics:phyx-mini-0735:phyxminiproblems-problemphyxmini0735-faultpoint, def:physics:phyx-mini-0735:phyxminiproblems-problemphyxmini0735-faultplanefigure}
  Evidence read from the primary raster and the scenario text. These fields record labels, line styles, qualitative motion, and the acute angle phi; they do not constrain the requested net magnitude
\end{definition}
\begin{proof}
  This is the declared model definition of Matches Primary Figure.
\end{proof}
\begin{definition}[Fault Component Laws]
  \label{def:physics:phyx-mini-0735:phyxminiproblems-problemphyxmini0735-faultcomponentlaws}
  \lean{PhyXMiniProblems.ProblemPhyXMini0735.FaultComponentLaws}
  \uses{def:physics:phyx-mini-0735:phyxminiproblems-problemphyxmini0735-displacementreadout, def:physics:phyx-mini-0735:phyxminiproblems-problemphyxmini0735-faultplanefigure}
  The component law for the free displacement vectors in the fault plane. It is stated in every length unit: AB is the vector sum of AC and AD, and the strike and dip directions are perpendicular
\end{definition}
\begin{proof}
  This is the declared model definition of Fault Component Laws.
\end{proof}
\begin{definition}[Matches Problem Data]
  \label{def:physics:phyx-mini-0735:phyxminiproblems-problemphyxmini0735-matchesproblemdata}
  \lean{PhyXMiniProblems.ProblemPhyXMini0735.MatchesProblemData}
  \uses{def:physics:phyx-mini-0735:phyxminiproblems-problemphyxmini0735-displacementmagnitudeinmeters, def:physics:phyx-mini-0735:phyxminiproblems-problemphyxmini0735-faultplanefigure}
  The two measured component magnitudes supplied in the question. No field here mentions the net magnitude or the recorded answer choice
\end{definition}
\begin{proof}
  This is the declared model definition of Matches Problem Data.
\end{proof}
% --- Archon declaration topology end ---
% --- Archon physics formalization source end ---
